\section{Evaluation}
\label{sec:evaluation}

Eight workflows arrive online on the recorded BLSC engine. Each node has two
A100-PCIe-40GB GPUs and serves one TP\,2 worker. Host-to-device copies use
the pinned rate measured on that partition, 23.98\,GB/s. The peer is a second
node of the same recorded type. Prefill is charged at 100\,$\mu$s per token
and decode at 2\,ms per token. Each workflow is one Agent. Its reads are a
captured sparse forward: NOSA-8B, NSA-8B, or MiniCPM4.1-8B on PG19 document
10146, 10321, or 10356. Arrivals are spaced at 0.45 of one workflow's service
time, so both workers stay busy. Restart keeps no checkpoint. DRAM-only keeps
a complete peer-DRAM checkpoint. \system adds the Top-$K$ blocks of the first
$L$ layers in peer HBM, with $L{=}4$, $K{=}16$, and a quota that starts at
16\,MiB. Router places the retry on the other node, KV Master admits the
checkpoint, and Adapter holds the quota.

With the failure removed, DRAM-only and \system finish every workflow at the
same time as Restart. The difference is 0\,ns on every healthy run in the
matrix: four prompt lengths on NOSA-8B document 10146, the same 8{,}192-token
arrival on NSA-8B and MiniCPM4.1-8B, and the other two NOSA documents.

The failure removes the worker running wf04 at 65\% of its invocation.
Figure~\ref{fig:eval-online} is the 24{,}576-token NOSA arrival. Workflows
released before the failure match the healthy run. wf04 finishes in
2{,}585.6\,ms with no failure, 4{,}102.9\,ms under \system, and
5{,}688.3\,ms under Restart. From wf04 through wf07, \system stays
1{,}585.4\,ms ahead of Restart. The queue before admission is 1{,}500.1\,ms
for both. Service is 2{,}459.6\,ms under Restart and 874.2\,ms under
\system. The checkpoint holds 15{,}872 of the 16{,}016 tokens already
produced, and the retry rebuilds the other 144. Restart rebuilds all
16{,}016. Router sends the retry from worker~1 to worker~0. The failed
worker stays down, so the two curves keep rising together and the checkpoint
remains a constant offset. DRAM-only lies within 0.217\,ms of \system on
every paired workflow, so the figure draws one line for both. \system places
64 HBM blocks. The quota is still 16\,MiB at the end of the run.

\begin{figure}[t]
  \centering
  \includegraphics[width=\columnwidth]{figures/eval-online-wct.pdf}
  \caption{Completion time for eight 24{,}576-token NOSA-8B workflows.
  The line marks the failure of wf04. DRAM-only finishes within 0.217\,ms
  of \system.}
  \label{fig:eval-online}
  \Description{Workflow completion before and after one mid-run failure.}
\end{figure}

\begin{center}
  \captionof{table}{Retry latency (ms) on NOSA-8B document 10146. Saved is
  Restart minus \system, and it is also the completion-time gap from wf04
  onward.}
  \label{tab:eval-retry}
  \Description{Retry latency of Restart and AegisKV at four prompt lengths.}
  \begin{tabular}{@{}rrrr@{}}
    \toprule
    Prompt & Restart & \system & Saved \\
    \midrule
    2{,}048 & 468.9 & 340.9 & 128.0 \\
    8{,}192 & 1{,}421.1 & 909.2 & 511.9 \\
    16{,}384 & 2{,}690.9 & 1{,}642.1 & 1{,}048.8 \\
    24{,}576 & 3{,}960.7 & 2{,}375.3 & 1{,}585.4 \\
    \bottomrule
  \end{tabular}
\end{center}

Figure~\ref{fig:eval-retry} and Table~\ref{tab:eval-retry} vary the prompt
length on the same document and the same failure. The saving grows from
128.0\,ms at 2{,}048 tokens to 1{,}585.4\,ms at 24{,}576 tokens. It tracks
the prefill of the checkpointed prefix. The queue grows with the prompt as
well, so retry latency falls by about two fifths rather than by the ratio of
tokens rebuilt. At 8{,}192 tokens, NSA-8B and MiniCPM4.1-8B, and NOSA
documents 10321 and 10356, all finish the retry within 6\,$\mu$s of the NOSA
document 10146 run. The captured selection changes which 64 blocks sit in
HBM. It does not move completion time beyond the 0.217\,ms DRAM gap.

\begin{figure}[t]
  \centering
  \includegraphics[width=\columnwidth]{figures/eval-retry-length.pdf}
  \caption{Retry latency of wf04 against prompt length. One GPU failure at
  65\% of the invocation. The vertical gap is the saving in
  Table~\ref{tab:eval-retry}.}
  \label{fig:eval-retry}
  \Description{Retry latency of Restart and AegisKV at four prompt lengths.}
\end{figure}

The single failure is also what the recorded GPU faults imply for a run of
this length. The InfiniteHBD trace covers 400 GPU servers and spans 345
days \cite{shou2025infinitehbd}. It contains 158 GPU faults on 97 servers.
The shortest GPU downtime is 267.8\,s and the median is 17.1\,h
(Figure~\ref{fig:eval-faults}). The longest arrival here, the failed
24{,}576-token run, ends 16.5\,s after the first release. Every recorded GPU
repair outlasts that arrival, so a fault that hits a live worker leaves the
worker down for the rest of the run. The trace is a pretraining-cluster log.
It supplies the downtime, not a serving request mix.

\begin{figure}[t]
  \centering
  \includegraphics[width=\columnwidth]{figures/eval-fault-downtime.pdf}
  \caption{GPU downtime in the InfiniteHBD trace. The line is the end of
  the longest serving run, 16.5\,s.}
  \label{fig:eval-faults}
  \Description{Cumulative distribution of recorded GPU fault durations.}
\end{figure}

The HBM copy stays small because $K{=}16$ covers little of these traces.
\label{sec:eval-native}
On the first four layers of NOSA, NSA, and MiniCPM, the deployed selector
hits 8.31\%, 4.20\%, and 4.06\% of distinct blocks. Recency matches the
first two and reaches 4.13\% on MiniCPM
(Figure~\ref{fig:eval-native-topk}).

\begin{figure}[t]
  \centering
  \includegraphics[width=\columnwidth]{figures/eval-native-topk.pdf}
  \caption{Distinct-block coverage on the three captured sparse models.
  Points average three documents; whiskers are the document range.}
  \label{fig:eval-native-topk}
  \Description{Coverage of the deployed selector, recency, and a future-demand oracle.}
\end{figure}
